% ======================================================================
% DOCUMENT PREAMBLE
% ======================================================================
\documentclass[11pt]{article}

% --- Code Listing Settings ---
\usepackage{listings}
\usepackage{xcolor}

\definecolor{codegreen}{rgb}{0,0.6,0}
\definecolor{codegray}{rgb}{0.5,0.5,0.5}
\definecolor{codepurple}{rgb}{0.58,0,0.82}
\definecolor{backcolour}{rgb}{0.95,0.95,0.92}

\lstdefinestyle{mystyle}{
    backgroundcolor=\color{backcolour},   
    commentstyle=\color{codegreen},
    keywordstyle=\color{magenta},
    numberstyle=\tiny\color{codegray},
    stringstyle=\color{codepurple},
    basicstyle=\ttfamily\footnotesize,
    breakatwhitespace=false,         
    breaklines=true,                 
    captionpos=b,                    
    keepspaces=true,                 
    numbers=left,                    
    numbersep=5pt,                  
    showspaces=false,                
    showstringspaces=false,
    showtabs=false,                  
    tabsize=2
}

\lstset{style=mystyle}

% --- Page Layout and Formatting ---
\usepackage[margin=1in]{geometry}
\usepackage{tabularx}
\usepackage{microtype}
\usepackage{booktabs}
\usepackage{float}
\usepackage[utf8]{inputenc}

% --- Mathematics and Physics Symbols ---
\usepackage{amsmath, amssymb, amsthm}
\usepackage{physics}
\usepackage{bm}
\usepackage{amsfonts}

% --- Graphics and Colors ---
\usepackage{graphicx}

% --- References and Hyperlinks ---
\usepackage{hyperref}
\hypersetup{colorlinks=true, linkcolor=blue, citecolor=blue, urlcolor=blue}
\usepackage[numbers]{natbib}
\usepackage{xurl}

% --- Theorem Environments ---
\newtheorem{definition}{Definition}
\newtheorem{assumption}{Assumption}
\newtheorem{postulate}{Postulate}
\newtheorem{proposition}{Proposition}

% ======================================================================
% TITLE AND AUTHOR INFORMATION
% ======================================================================
\title{The Cosmological Lensing Effect:\\ Geometric Curvature as an Alternative to Accelerated Expansion}
\author{Binyamin Tsadik Bair-Mosheh\thanks{Head of Physics at System Advanced Laboratories Ltd., Yavne, Israel. Email: binyamin.tsadik@gmail.com. ORCID: 0000-0003-4329-7740 }}
\date{\today}

% --- Hyphenation and Emergency Settings ---
\emergencystretch=2em
\hyphenation{Pantheon supercluster}

% ======================================================================
% MAIN DOCUMENT CONTENT
% ======================================================================
\begin{document}
\maketitle

\begin{abstract}
The inference of an accelerating universe relies on the interpretation of Type Ia supernovae within the standard Friedmann-Lema\^itre-Robertson-Walker metric. This metric retains the weak-field approximation by treating the macroscopic background as a smooth, expanding coordinate grid. In this work, we investigate the consequences of lifting this assumption and treating the background gradient as an evolving refractive medium. By applying the 1923 Gordon optical metric to a curved spacetime, we derive the longitudinal lensing effects driven by the background curvature, quantified by the scalar velocity field $v_0$. To remove prior assumptions and isolate the kinematics for empirical curve fitting, we apply an analytic inversion to the Pantheon+ supernova dataset, transforming the raw observational phase space ($z, \mu$) directly into de-lensed physical coordinates ($r_{true}, v_0$). By exploring this cleaned data without the weak-field constraint, our analysis suggests that continuous optical path lengthening significantly reduces the magnitude anomaly attributed to dark energy. This is accomplished by using a single free parameter $K$ which converges on a value that is equivalent to the Hubble radius $R_h$, effectively becoming a zero-free-parameter model. Evaluating the exact metric refractive index $n(r) = (1 - v_0^2/c^2)^{-1}$, the model fits the Pantheon+ supernova sample competitively with flat $\Lambda$CDM ($\Delta\chi^2 = +26$ vs.\ 1-parameter Planck $\Lambda$CDM, while improving upon pure coasting FLRW by $\Delta\chi^2 = -64$), and reproduces the transverse BAO width ($D_M/r_d$) of DESI Year 1 across $z \le 1.5$ ($\chi^2_{DM} = 24.7$ vs.\ $36.1$ for Pantheon+-calibrated $\Lambda$CDM). However, the 1-parameter pure-vacuum continuum does not capture DESI's intermediate-redshift line-of-sight $D_H(z)$ inflection or the Ly$\alpha$ ($z=2.33$) residual, which is consistent with the omission of baryonic matter ($\Omega_b \approx 0.049$) from the pure-vacuum equation of state. Extending the model to a full multi-epoch BAO solution is left to future work.
\end{abstract}

\vspace{0.5cm}
\noindent \textbf{Keywords:} Cosmological Lensing, Refractive Metric, Dark Energy, Hubble Tension, Gravitational Time Dilation

% ======================================================================
% Section 1: Introduction: Measurement and Interpretation
% ======================================================================
\section{Introduction: Measurement and Interpretation}
\label{sec:macroscopic_metric}

Cosmological observation is fundamentally constrained by the nature of signal propagation along the past light cone ($ds^2 = 0$). Consequently, a photon observed today from a high-redshift source has not merely traveled a geometric distance; it has traversed a continuous gradient in the background metric. 

The assumption of a flat, unperturbed background originated with Minkowski's formalization of the four-dimensional spacetime continuum in 1908~\cite{Minkowski1909}. While Einstein fundamentally dismantled the rigidity of this background in the formulation of General Relativity, he explicitly reintroduced the flat background metric ($\eta_{\mu\nu}$) as a mathematical simplification. This weak-field approximation was utilized to linearize the field equations and demonstrate their reduction to Newtonian mechanics~\cite{Einstein1916}. 

In standard cosmological frameworks, this approximation has subtly persisted, often treating macroscopic space as a smooth, expanding coordinate grid overlaid with localized gravitational noise. The primary objective of this paper is to eliminate this historical simplification. Rather than treating the background metric as a passive arena, we treat the macroscopic gradient as a relativistic refractive medium.

\subsection{The Effective Refractive Index of the Vacuum}
Following standard treatments of the effective optical metric and null congruences in curved spacetime~\cite{Gordon1923, Poisson2004}, we posit a direct correspondence between the geometry of the null geodesic and the optical properties of the vacuum. 

In a spherically symmetric gravitational potential $\Phi(r)$ governed by the reciprocal compliance metric ($g_{00} = -(1 + 2\Phi/c^2)$ and $g_{rr} = (1 + 2\Phi/c^2)^{-1}$), recovering the full relativistic time delay and deflection angle~\cite{Einstein1916, Eddington1920} rather than the classical S\"oldner half-deflection~\cite{Soldner1801} requires incorporating both the temporal lapse and the spatial metric component~\cite{Poisson2014}. Solving the null condition ($ds^2 = 0$) yields the exact effective refractive index $n(r) \equiv c / (dr/dt)$:
\begin{equation}
\label{eq:n_weak}
n(r) = \left(1 + \frac{2\Phi(r)}{c^2}\right)^{-1}
\end{equation}

Recognizing that the classical escape velocity is defined by $v_{\mathrm{esc}}^2 = -2\Phi$ (for an attractive potential vanishing at infinity), we establish the exact constitutive relation for the refractive vacuum:
\begin{equation}
\label{eq:n_vesc}
n(r) = \left(1 - \frac{v_{\mathrm{esc}}^2(r)}{c^2}\right)^{-1}
\end{equation}
The local gravitational acceleration is thus intrinsically linked to the spatial gradient of this index, physically scaling the rate of refractive transport.

\subsection{The Effective Static Mapping of an Evolving Background}
The universe has been consistently evolving from a denser past to its current state of relative sparsity. In standard cosmology, this evolution is treated temporally via a time-dependent scale factor $a(t)$. For this analysis, we adopt the background kinematic frame $R_h = c \cdot t$ (linear horizon expansion), which is derived formally from an elastic $w = -1/3$ vacuum equation of state in our companion formalization~\cite{RefractivePackage2026}. While the refractive optical effects derived in Sections~\ref{sec:gravitational_gradient} and~\ref{sec:results_discussion} follow directly from the metric potential, the specific redshift-distance relation integrated in Appendix~\ref{app:data_reconstruction} leverages this linear kinematic baseline.

Cosmological observation is fundamentally constrained to the past light cone ($ds^2 = 0$). Because radial distance and look-back time are coupled for propagating photons, the temporal evolution of the cosmic medium is observationally isomorphic to a spatial gradient. When we observe a high-redshift source, we are looking deeply into this evolution to a higher curvature and denser past. For a null geodesic in a flat FLRW metric, the coordinate velocity is governed by the scale factor:
\begin{equation}
ds^2 = 0 \implies v_c = \left| \frac{dr}{dt} \right| = \frac{c}{a(t)}
\end{equation}

Conversely, in a static optical metric, the coordinate velocity of light is governed by the refractive index:
\begin{equation}
v_c = \left| \frac{dr}{dt} \right| = \frac{c}{n(r)}
\end{equation}

Equating these kinematic constraints dictates that along the null geodesic, the temporal evolution acts observationally as a spatial refractive index. Rather than modeling photon propagation through a flat, temporally expanding coordinate grid, we can map this continuous change in background density directly onto an effective spatial potential, $\Phi(r)$. Under this equivalence, the cosmological redshift obeys the same functional form as relativistic time dilation:
\begin{equation}
1+z = \frac{a(t_0)}{a(t_e)} \equiv \sqrt{\frac{g_{00}(r_0)}{g_{00}(r_e)}}
\end{equation}

Therefore, the continuous spatial-temporal dilution of the continuum can be encoded geometrically by modeling it as a Gradient-Index Lens (GRIN). The apparent "expansion rate" is thus reinterpreted, at least partially, as the refractive transport of signals through the potential basin of an earlier universe.

\subsection{Transverse vs. Longitudinal Lensing}
Consistent with the Gordon metric formulation, the trajectory of light in this potential satisfies Fermat's Principle of least time ($\delta \int n \, ds = 0$)~\cite{BornWolf}. While the application of this principle is universally accepted, its standard usage in cosmology is currently limited to the transverse plane.

\paragraph{The Transverse Constraint (Standard Lensing):}
When light grazes a local mass distribution, the transverse component of the gradient ($\nabla_{\perp} n$) generates the standard gravitational deflection angle: $\alpha \simeq (2/c^2) \int \nabla_{\perp} \Phi \, ds$. Because these effects are well-explored~\cite{Bartelmann2001}, standard cosmology treats lensing as localized "noise" against a smooth expanding background.

\paragraph{The Longitudinal Constraint (The Macroscopic Monopole):}
Standard cosmological lensing models the influence of localized intervening mass via the transverse gradient ($\nabla_{\perp} n$). As first demonstrated by Weinberg (1976)~\cite{Weinberg1976}, photon flux conservation dictates that these transverse magnifications and demagnifications average to unity across the sky. Consequently, stochastic weak lensing contributes only a negligible, sub-percent scatter to the distance modulus in large supernovae samples.  

However, the primary objective of this work is applying the optical metric \textit{longitudinally} ($\nabla_{\parallel} n$) along the entire line of sight. Because the photon is continuously traveling out of the potential well of the early universe, this gradient acts as a global, non-canceling monopole. Under this geometry, Fermat's principle dictates that the photon experiences a continuous Shapiro Delay~\cite{Shapiro1964}. This accumulation of optical path length forces the apparent observation distance to diverge from the true geometric distance:
\begin{equation}\label{eq:longitudinal_delay}
d_{\mathrm{optical}} = \int_{0}^{d_{\mathrm{geo}}} n(r) \, dr > d_{\mathrm{geo}}
\end{equation}
Because $n(r) > 1$ continuously inside this evolving potential well, the optical path is necessarily longer than the geometric path. Unlike stochastic weak lensing, this longitudinal delay cannot average to zero across different lines of sight; it creates a systematic, isotropic baseline shift. We propose that the apparent "acceleration" of the universe (Dark Energy) is at least partially attributable to the observational consequence of this uncorrected monopole lens.

% ----------------------------------------------------------------------
% Section 2: The Gravitational Gradient: Derivation of the Potential Effects
% ----------------------------------------------------------------------
\section{The Gravitational Gradient: Derivation of the Potential Effects}
\label{sec:gravitational_gradient}

In order to de-lens the supernova data, we must first quantify the effects that an evolving background has on the redshift signal. Traditionally, cosmology acknowledges that the universe is expanding and that this expansion must be accompanied by a reduction in average baryon density; however, numerical approximations typically treat photon propagation as immune to the evolving background by assuming a background metric where $g_{00} \approx -1$ everywhere. Here, we rigorously derive the optical and relativistic consequences of retaining the macroscopic potential gradient $\nabla \Phi$.

We define the potential gradient of the line of sight by the characteristic velocity of the potential well, $v_0(r)$. This single scalar function governs four distinct, coupled physical effects: gravitational time dilation, kinematic velocity coupling, kinematic time dilation, and refractive lensing.

\subsection{The Potential Formalism}
We begin with the reciprocal compliance metric for a spherically symmetric potential $\Phi(r)$, where the spatial packing scales as the inverse of the temporal lapse ($g_{rr} = -g_{00}^{-1}$). Defining $v_0(r)$ as the classical escape velocity associated with the potential depth relative to the observer ($v_0^2(r) \equiv -2\Phi(r)$), the line element is:
\begin{equation}
\label{eq:v0_metric}
ds^2 = -\left(1 - \frac{v_0^2(r)}{c^2}\right)c^2 dt^2 + \left(1 - \frac{v_0^2(r)}{c^2}\right)^{-1} dr^2.
\end{equation}

\subsection{Effect 1: Gravitational Time Dilation}
The metric term $g_{00}$ governs the flow of proper time $\tau$. For a source comoving with the potential at emission depth $r_{source}$, the relationship between proper time and coordinate observer time $t$ is $d\tau = dt \sqrt{-g_{00}}$. 

Physical processes deep in the potential well appear dilated to the observer by the factor:
\begin{equation}
\label{eq:grav_dilation}
\gamma_{\Phi}(v_0) = \frac{1}{\sqrt{1 - v_0^2/c^2}}.
\end{equation}
Because the frequency of an emitted photon is the inverse of the temporal period, the metric dilation causes the photon's observed frequency to drop. This yields the classical Gravitational Redshift ($1 + z_{\Phi} = \gamma_{\Phi}$), historically verified by terrestrial experiments~\cite{Pound1960}.

\subsection{Effect 2: Kinematic Velocity Coupling}
Adopting the assumption that massive bodies are in a state of free-fall, the recession velocity $v_{kin}$ of a galaxy equates to the local escape velocity of the potential well:
\begin{equation}
\label{eq:kinematic_coupling}
v_{kin}^2(r) \approx v_0^2(r) = -2\Phi(r).
\end{equation}
This macroscopic motion stretches the emitted wavelength, producing a classical kinematic Doppler shift: $1 + z_{class} = 1 + v_{kin}/c = 1 + v_0/c$.

\subsection{Effect 3: Kinematic Time Dilation}
Because the source galaxy is in physical motion relative to the observer ($v_{kin}$), it is also subject to the time dilation effects of Special Relativity, $\gamma_{kin}$. Through the velocity coupling established in Eq.~\ref{eq:kinematic_coupling}, this kinematic dilation is numerically identical to the gravitational dilation, such that $\gamma_{kin}(v_{kin}) = \gamma_{\Phi}(v_0)$.

The total observed redshift of the photon ($z_{tot}$) is the product of the gravitational term (Effect 1) and the full kinematic Doppler term (which combines the classical spatial stretch of Effect 2 with the kinematic dilation of Effect 3), all accounted for by a single scalar function $v_0(r)$:
\begin{equation}
1 + z_{tot} = (1 + z_{\Phi}) (1 + z_{kin}) \approx \gamma_{\Phi}(v_0) \times \left[ \gamma_{kin}(v_0) \left(1 + \frac{v_0}{c}\right) \right].
\end{equation}
Because $\gamma_{kin} = \gamma_{\Phi}$, the compound dilation simplifies to $\gamma^2(1+v_0/c)$, which forms the mathematical basis for the analytic data inversion utilized in Appendix~\ref{app:data_reconstruction}.

\subsection{Effect 4: Refractive Index}
The spatial curvature of the potential governs the coordinate speed of light $v_c(r) = dr/dt$. Solving for the radial null geodesic ($ds^2=0$) in Eq.~\ref{eq:v0_metric}:
\begin{equation}
\left(1 - \frac{v_0^2}{c^2}\right)^{-1} \left(\frac{dr}{dt}\right)^2 = \left(1 - \frac{v_0^2}{c^2}\right) c^2 \implies \frac{dr}{dt} = c\left(1 - \frac{v_0^2}{c^2}\right).
\end{equation}
We define the exact refractive index $n(r) \equiv c / (dr/dt)$ directly from the radial null geodesic (matching Eq.~\ref{eq:n_vesc}) as:
\begin{equation}
n(r) = \left(1 - \frac{v_0^2(r)}{c^2}\right)^{-1}.
\end{equation}
Because any observed source lies at a finite redshift $z < \infty$ ($v_0(z) = c z/(1+z) < c$), this exact refractive index remains strictly finite along the entire observed past light cone while naturally diverging ($n \to \infty$) as $v_0 \to c$ at the causal horizon boundary.

% ======================================================================
% Section 3: Results and Discussion
% ======================================================================
\section{Results and Discussion}
\label{sec:results_discussion}

By applying the potential gradient formalism derived in Section~\ref{sec:gravitational_gradient}, we identify a critical optical correction that has been traditionally overlooked.

\subsection{The Velocity Potential Profile}
We applied a direct analytic inversion to the Pantheon+ supernova dataset~\cite{PantheonPlus, SHOES2022} using the iterative rectification pipeline detailed in Appendix~\ref{app:data_reconstruction}. Unlike traditional analyses that fit model parameters ($\Omega_m, \Omega_\Lambda$) to the data, we transformed each individual supernova observation from redshift-magnitude space ($z, \mu$) into de-lensed physical coordinates ($v_0, r_{true}$).

This transformation yields a smooth, coherent velocity profile $v_0(r)$ that exhibits clear saturation at high distances. This profile emerges naturally from the data without the need for smoothing, binning, or the imposition of $\Lambda$CDM priors.

\subsection{The Geometric Shape of the Universe}
To quantify this empirically derived profile, we fitted the data to two candidate distributions (Appendix~\ref{app:potential_law}):
\begin{enumerate}
    \item \textbf{Exponential Fit:} A thermodynamic saturation distribution ($v \propto 1 - e^{-r/K}$) with a converged scale radius $K \approx 4200$ Mpc ($4200.4$ Mpc).
    \item \textbf{Sigmoidal Fit:} A geometric saturation distribution with scale radius $K \approx 3029$ Mpc and shape parameter $\eta \approx 1.18$.
\end{enumerate}

Both models achieved high statistical significance ($R^2 > 0.99$), suggesting that the local Hubble distribution ($v \propto r$) is a linear approximation of the curved potential that saturates ($v_0 \to c$) at the cosmic horizon. Notably, the Exponential fit recovers a local gradient of $H_0 = c/K \approx 71.4$ km/s/Mpc, placing it squarely within the measurement window of the current Hubble tension.

\subsection{Refractive Dimming as Dark Energy}
Cosmology currently attributes the anomalous dimming of high-redshift supernovae to an accelerating metric, driven by "Dark Energy." However, this interpretation fundamentally relies on the mathematical crutch of an optically empty, flat background metric ($n=1$).

By evaluating the exact refractive index $n(r) = (1 - v_0^2/c^2)^{-1}$ in terms of background curvature (Appendix~\ref{app:vacuum_metric}), we quantify a systematic optical path lengthening. Because $n(r) > 1$ continuously inside the potential well, the apparent distance light travels ($r_{opt}$) necessarily diverges from the true geometric coordinate ($r_{true}$) via the continuous optical integral: $r_{opt} = \int n(r)\,dr$.

As derived in Appendix~\ref{app:dark_energy_origin}, the magnitude anomaly ($\Delta \mu_{DE}$) is not a measurement of metric acceleration, but the signature of this uncorrected refractive lens:
\begin{equation}
    \Delta \mu_{DE} = 5 \log_{10} \left( \frac{r_{opt}}{r_{true}} \right)
\end{equation}
\textit{When this optical correction is applied to the Pantheon+ dataset, the "accelerating" signal flattens to zero across all redshifts.} This demonstrates that for the supernova distance-redshift relation, the dimming can be accounted for by the continuous longitudinal Shapiro delay of the macroscopic potential well. \textbf{This suggests that the Dark Energy signature in SNe Ia can be modeled geometrically as an artifact of longitudinal lensing.}

\subsection{Resolution of the Hubble Tension}
\label{sec:Hubble_Tension}
The discrepancy between local ($H_0 \approx 73$ km/s/Mpc) and cosmic ($H_0 \approx 67.4$ km/s/Mpc) measurements is resolved by the non-linear saturation of the velocity profile. Equating the exponential profile $v_0(r) = c(1 - e^{-r/K})$ to the redshift inversion $v_0(z) = c z/(1+z)$ (Appendix~\ref{app:data_reconstruction}) yields $e^{-r/K} = (1+z)^{-1}$, or $r_{true}(z) = K \ln(1+z)$. Taking the radial derivative of $v_0(r)$ shows that the local velocity-distance slope ($H_{slope} \equiv dv_0/dr$) decays with distance and redshift:
\begin{equation}
    H_{slope}(z) \equiv \frac{dv_0}{dr} = \frac{c}{K}e^{-r/K} = \frac{c}{K(1+z)} = \frac{H_0}{1+z}
\end{equation}
Crucially, the local velocity-distance slope $H_{slope}(z) = dv_0/dr$ and the BAO line-of-sight expansion rate $H_{BAO}(z) \equiv c / (dr_{true}/dz) = (c/K)(1+z) = H_0(1+z)$ (Section~\ref{sec:bao_delensing}) are distinct observables related by $H_{BAO}(z)/H_{slope}(z) = (1+z)^2$, which coincide at $z = 0$ ($H_{slope}(0) = H_{BAO}(0) = H_0 = c/K \approx 71.4$ km/s/Mpc).

Consequently, the Hubble parameter measured from local velocity-distance gradients (SH0ES) operates at $z \approx 0$, recovering the maximum tangent slope $H_0 = c/K \approx 71.4$ km/s/Mpc. Conversely, global fits that force a single baseline across the full depth of the potential well infer a lower effective value. The "tension" is therefore not a measurement error, but a predictable consequence of evaluating a saturating exponential potential across different redshift baselines.

\subsection{The Age Illusion and JWST Anomalies}
When evaluating the evolutionary time at high redshifts ($z=14$), standard matter-dominated expansion baselines predict that galaxies are constrained by a narrow window (affording only $\approx 137$ Myr between $z=20$ and $z=14$). Our relativistic analysis (Appendix~\ref{app:time_dilation}) reveals that the continuous temporal scaling of the spatial metric dictates a significantly extended chronological boundary.

Using the chronometer derived from our geometric scale radius ($K \approx 4200.4$ Mpc), the proper time ($\Delta t_{evol}$) available for structure formation between "first light" ($z=20$) and the epoch of observation ($z=14$) is:
\begin{equation}
    \Delta t_{evol} \approx 261 \text{ Myr} \quad (\text{Exponential Distribution})
\end{equation}
This effectively doubles ($\approx 1.9\times$) the early temporal envelope available for stellar evolution compared to an unsaturated expansion baseline, accommodating the presence of the massive, mature galaxies observed by JWST~\cite{Labbe2023}.

\subsection{Proposed Verification via BAO Delensing}
\label{sec:bao_delensing}
The geometric predictions of this refractive effect can be independently tested by analyzing the large-scale clustering of Baryon Acoustic Oscillations (BAO). Modern spectroscopic surveys, such as SDSS DR12, eBOSS DR16, and DESI Year 1~\cite{Alam2021, Adame2024}, map the primordial sound horizon using distinct geometric observables. Because the Cosmological Lensing Effect predicts specific, diverging distortions for longitudinal versus transverse optics, BAO clustering provides an ideal methodology for independent verification.

If the universe acts as a refractive lens, applying the derived velocity field (adopting the fit of $K \approx 4200$ Mpc) to the raw BAO observables yields specific baseline predictions across these spatial dimensions:
\begin{itemize}
    \item \textbf{Longitudinal Depth ($D_H$):} The true physical line-of-sight depth of the BAO ripples matches the spatial derivative of the logarithmic true coordinate geometry ($r_{true} = K \ln(1+z)$), yielding $D_{H,true}(z) = \frac{d}{dz}(K \ln(1+z)) = \frac{K}{1+z}$.
    \item \textbf{Transverse Width ($D_M$):} The apparent angular width of the sound horizon is subject to continuous magnification, tracking the predicted apparent distance of the gradient-index lens. Substituting $1 + z' = e^{r'/K}$ ($dr' = \frac{K}{1+z'}dz'$) and the exact refractive index $n(z') = \left(1 - \frac{v_0^2}{c^2}\right)^{-1} = \frac{(1+z')^2}{2z'+1}$ into the line-of-sight integral ($r_{opt} = \int_{0}^{r_{true}} n(r')\,dr'$) yields a closed-form analytic function of redshift:
    \begin{equation}
        r_{opt}(z) = K \int_0^z \frac{1+z'}{2z'+1}\,dz' = K \left[ \frac{z}{2} + \frac{1}{4}\ln(2z+1) \right]
    \end{equation}
    \item \textbf{Volumetric Distance ($D_V$):} For low-redshift surveys (e.g., 6dFGS) where transverse and longitudinal dimensions are blended into a single volume-averaged measure, the apparent geometry converges toward the true physical volume, geometrically constrained by $D_{V,true} = (z \cdot r_{true}^2 \cdot dr/dz)^{1/3}$.
    \item \textbf{Restoration of Sphericity:} The Alcock-Paczy\'nski (AP) test measures the ratio between the angular size and redshift depth of the sound horizon, $F_{AP}(z) = D_M(z)/D_H(z) = (1+z)[z/2 + \frac{1}{4}\ln(2z+1)]$, which is completely independent of both $K$ and $r_d$.
\end{itemize}

Evaluating these baseline predictions against the DESI Year 1 dataset~\cite{Adame2024} clarifies both the reach and the physical boundary of the 1-parameter pure-vacuum potential. Across the galaxy and quasar tracers ($z \le 1.5$), the exact optical integral $r_{opt}(z)$ reproduces DESI's transverse $D_M(z)/r_d$ measurements with $\chi^2_{DM} = 24.7$ (under the Planck acoustic ratio $K/r_d = 30.24$, outperforming Pantheon+-calibrated flat $\Lambda$CDM at $\chi^2_{DM} = 36.1$), while the parameter-free Alcock-Paczy\'nski ratio matches DESI at $z = 0.51$ ($F_{AP} = 0.650$ vs.\ $0.649 \pm 0.022$, $+0.06\sigma$). However, because the 1-parameter curve $v_0(r) = c(1 - e^{-r/K})$ models only the macroscopic vacuum continuum ($w = -1/3$, scaling as $(1+z)^2$) and omits the $\sim 5\%$ baryonic matter density ($\Omega_b \approx 0.049$, which dilutes volumetrically as $(1+z)^3$), it does not capture the intermediate-redshift ($z \sim 0.7\text{--}1.3$) inflection in $D_H(z)$ or the high-redshift Ly$\alpha$ ($z = 2.33$) residual. Extending the model to a full multi-epoch BAO solution therefore requires coupling the baryonic matter density to the vacuum potential within the covariant conservation law, which we leave to dedicated future work.

\subsection{Independent Verification via CMB Schlieren Imaging}
While this work establishes the Cosmological Lensing Effect through the longitudinal lensing of supernovae, the geometric signature of the refractive vacuum must also be detectable via independent macroscopic methods. We propose that the uniform Cosmic Microwave Background (CMB) backlight at $z \approx 1100$ effectively transforms the observable universe into a macroscopic Schlieren interferometer~\cite{Toepler1864}.

As detailed in Appendix~\ref{app:cosmic_schlieren}, standard cosmological analysis typically subtracts low-multipole lensing potentials as systemic noise. However, the spatial curvature of the vacuum dictates a global convergence ($\kappa_{global} \propto \nabla^2 n$). This is not statistical noise; it is a structural mandate of the metric that generates a systematic, isotropic magnification of the CMB power spectrum.

We hypothesize that the refractive continuum may provide a physical mechanism for the persistent $A_L$ lensing amplitude anomaly observed by Planck~\cite{Planck2018} (where $A_L \approx 1.18$). If this anomaly is not a statistical fluke, it may represent the literal Schlieren signature of CMB photons traversing the uncorrected global refractive gradient. While deriving the exact scalar amplitude of $A_L$ from $K$ is left to future work, tracking this directional structural variance provides a potential independent methodology to map the 3D topology of the macroscopic vacuum.

\subsection{Conclusion: The Cosmological Lensing Effect}
By restoring the refractive geometry inherent to classical General Relativity, the magnitude anomaly in Type Ia supernovae attributed to dark energy can be modeled geometrically without invoking metric acceleration.

This suggests that several prevailing anomalies of modern cosmology (including the magnitude dimming of supernovae, the Hubble tension, and the massive early galaxies observed by JWST) can be treated as predictable, systematic artifacts generated by evaluating a gradient-index optical medium as if it were a flat, empty coordinate grid. 

Additionally, this optical framework is rigorously falsifiable and invites independent confirmation. By treating the universe as a continuous macroscopic lens, the geometry of the vacuum can be tested across multiple cosmological probes: de-lensing the transverse and longitudinal dimensions of Baryon Acoustic Oscillations (where the 1-parameter pure-vacuum baseline succeeds for $z \le 1.5$ transverse $D_M/r_d$ while isolating the baryonic matter contribution required at higher redshifts), and evaluating the uniform Cosmic Microwave Background as a global Schlieren interferometer.

We conclude that we are potentially observing the cosmos through the uncorrected lens of its own gravity. Recognizing this \textit{Cosmological Lensing Effect} strips away the optical distortions of the macroscopic vacuum, allowing us to approach a more physical model of the universe.

% ======================================================================
% STATEMENTS AND DECLARATIONS
% ======================================================================
\section*{Statements and Declarations}

\textbf{Funding:} The author did not receive support from any organization for the submitted work. No funding was received to assist with the preparation of this manuscript.

\textbf{Competing Interests:} The author has no relevant financial or non-financial interests to disclose.

\textbf{Data Availability:} The observational data underlying this article were derived from sources in the public domain: the Pantheon+ Type Ia supernova compilation (\url{https://github.com/PantheonPlusSH0ES/DataRelease}) and the Planck 2018 Cosmological Parameters (\url{https://pla.esac.esa.int}). The derived data generated in this research, along with the Python code used to perform the analytic data inversion and physical model fitting, are available as an Online Resource~\cite{RefractivePackage2026}.

% ======================================================================
% APPENDICES
% ======================================================================
\appendix

\section{Appendix A: Derivation of the Velocity Potential Data}
\label{app:data_reconstruction}

To investigate the shape of the cosmic potential well without imposing a specific cosmological model (such as $\Lambda$CDM), we transformed the observational data from the Pantheon+ Supernova dataset directly into physical coordinates.

Two distinct datasets were generated to distinguish between the \textit{apparent} geometry (distorted by curvature) and the \textit{physical} geometry (the underlying matter distribution):
\begin{enumerate}
    \item \textbf{Dataset 1 (Apparent):} The raw coordinate inversion assuming geometric optics ($d_{opt} \approx r$). This represents the universe as seen through the lens.
    \item \textbf{Dataset 2 (Rectified):} The de-lensed coordinates ($r_{true}$) obtained by subtracting the refractive path delay. This represents the physical structure of the potential well.
\end{enumerate}

\subsection*{Data Selection and Pre-Processing}
We utilized the \texttt{Pantheon+SH0ES.dat} compilation, selecting 1701 Type Ia supernovae with high-quality distance modulus measurements ($z_{obs} > 0.001$). The raw observables for each object are the cosmological redshift $z_{obs}$ and the calibrated distance modulus $\mu_{obs}$.

\subsection{Step 1: Analytic Inversion of the Energy State ($v_0$)}

We define $v_0$ as the escape velocity of the vacuum relative to the observer's location (where $v_0=0$). To determine this value from the observed redshift, we account for the two distinct relativistic effects acting on the photon:

\begin{enumerate}
    \item \textbf{Kinematic Doppler Shift ($z_k$):} Caused by the recession velocity of the local frame ($v_{flow}$).
    \begin{equation}
        1 + z_k = \sqrt{\frac{1+v_{flow}/c}{1-v_{flow}/c}} = \gamma(v_{flow}) \left(1 + \frac{v_{flow}}{c}\right)
    \end{equation}
    \item \textbf{Gravitational Redshift ($z_g$):} Caused by the photon climbing out of the potential well depth $\Phi = -v_0^2/2$.
    \begin{equation}
        1 + z_g = \frac{1}{\sqrt{1 - v_0^2/c^2}} = \gamma(v_0)
    \end{equation}
\end{enumerate}

We assume a critical "Free Fall" condition where the kinematic flow velocity equals the potential escape velocity ($v_{flow} = v_0$). The total observed redshift is the product of these two factors:
\begin{equation}
    1 + z_{tot} = (1 + z_k)(1 + z_g) = \left[ \gamma \left(1 + \frac{v_0}{c}\right) \right] \times \gamma
\end{equation}
\begin{equation}
    1 + z_{tot} = \gamma^2 \left(1 + \frac{v_0}{c}\right) = \left( \frac{1}{1 - v_0^2/c^2} \right) \left(1 + \frac{v_0}{c}\right)
\end{equation}

Expanding the denominator as a difference of squares ($1 - v^2/c^2 = (1-v/c)(1+v/c)$) reveals the algebraic cancellation:
\begin{equation}
    1 + z_{tot} = \frac{1 + v_0/c}{(1 - v_0/c)(1 + v_0/c)} = \frac{1}{1 - v_0/c}
\end{equation}

Solving for $v_0$ yields the analytic inversion:
\begin{equation}
    v_0 = c \cdot \left( \frac{z_{tot}}{1 + z_{tot}} \right)
\end{equation}
This establishes the energy state of each object solely from its redshift, independent of its spatial distance.

\subsection{Step 2: Recovery of the Apparent Coordinate ($r_{opt}$)}
To determine the radial coordinate we invoke the Etherington distance-duality relation, $d_L = (1+z)^2 d_A$, which holds in any metric theory of gravity where photons follow null geodesics and photon number is conserved, including the static refractive metric of Eq.~\ref{eq:v0_metric}. Calculating the luminosity distance $d_L$ from the observed magnitude $\mu_{obs}$, we define the apparent optical distance $r_{opt} \equiv (1+z)d_A$ by removing one redshift factor $(1+z)$:
\begin{equation}
    r_{opt} = \frac{d_L}{1 + z} = \frac{10^{(\mu_{obs} - 25)/5}}{1 + z}
\end{equation}
The dataset $(r_{opt}, v_0)$ constitutes \textbf{Dataset 1}. It represents the apparent optical profile of the universe including the longitudinal refractive delay caused by the intervening potential.

\subsection{Step 3: Iterative Optical Rectification (Dataset 2)}
The apparent distance $r_{opt}$ is an overestimate of the true depth because the exact refractive index of the potential, $n(r) = (1 - v_0^2/c^2)^{-1}$, lengthens the optical path. To recover the physical coordinate distance $r_{true}$, we must "de-lens" the data by solving the optical path integral equation:
\begin{equation}
    r_{opt} = \int_0^{r_{true}} n(r') \, dr' = \int_0^{r_{true}} \left( 1 - \frac{v_0(r')^2}{c^2} \right)^{-1} dr'
\end{equation}

Since the refractive index $n(r)$ depends on the potential profile $v_0(r)$, which is itself a function of the unknown true distance $r_{true}$, this creates a circular dependency. We resolved this via an iterative self-consistency loop:
\begin{enumerate}
    \item We initialized the potential with the apparent fit parameters ($K_{init} \approx 4400$ Mpc).
    \item We numerically inverted the optical integral to derive a candidate set of physical distances $r_{true}$.
    \item We fitted the Exponential model ($v_0 = c(1-e^{-r/K})$), as devised in Appendix~\ref{app:potential_law}, to these new coordinates to derive an updated scale radius $K_{new}$.
    \item We repeated steps 2 and 3 until the parameter $K$ converged.
\end{enumerate}

The process converged rapidly (within 4 iterations) to a stable physical scale radius of \textbf{$K \approx 4200.4$ Mpc}.

\subsection{Output Dataset}
The result of this pipeline is \textbf{Dataset 2} $(r_{true}, v_0)$, which represents the \textit{Rectified Potential Profile}. This clean dataset serves as the input for the final curve-fitting analysis presented in Appendix~\ref{app:potential_law}. It demonstrates that the "acceleration" signal (Dark Energy) in the supernova Hubble diagram is contained between the apparent optical path $r_{opt}$ (Step 2) and the physical coordinate depth $r_{true}$ (Step 3).

% ----------------------------------------------------------------------
% Appendix B: Analytical Determination of the Velocity Potential
% ----------------------------------------------------------------------
\section{Appendix B: Analytical Determination of the Velocity Profile}
\label{app:potential_law}

In Appendix~\ref{app:data_reconstruction}, we reconstructed the escape velocity profile $v_0(r)$ directly from the observational data, creating a clean dataset of physical coordinates $(r_{\text{true}}, v_0)$ independent of model priors. In this appendix, we analyze this profile to identify the underlying "best fit" to the background potential.

\begin{figure}[H]
    \centering
    \includegraphics[width=0.85\textwidth]{velocity_potential_fit.png}
    \caption{\textbf{The Saturated Potential of the Universe.} The reconstructed escape velocity profile ($v_0$ vs $r_{true}$) for 1701 Type Ia supernovae (grey points). The data clearly deviates from a linear trajectory (unbounded expansion) and saturates at the horizon ($v_0 \to c$). The Exponential fit (blue, $K = 4200$ Mpc) models this saturation with high precision ($R^2 = 0.9929$), revealing the finite depth of the cosmic potential well.}
    \label{fig:velocity_potential}
\end{figure}

\subsection{Candidate Fits}

We tested two candidate functions that satisfy the horizon boundary condition ($v_0 \to c$ as $r \to \infty$).

\textbf{Fit A: Exponential Saturation}
This is the simplest one-parameter model for asymptotic bounded growth. It is distinct in that it maintains a linear behavior at the origin ($v_0 \propto r$) recovering the local Hubble Law, before smoothly saturating at the horizon. It is characterized by a single scale radius $K_{exp}$:
\begin{equation}
    v_0(r) = c \cdot \left( 1 - e^{-r/K_{exp}} \right)
\end{equation}

\textbf{Fit B: Sigmoidal Saturation}
This two-parameter model introduces a shape parameter $\eta$ to test for a non-linear potential core. Unlike the linear origin of the exponential fit, a shape parameter of $\eta > 1$ creates an S-curve, suppressing the local velocity gradient before it accelerates and eventually saturates at the horizon:
\begin{equation}
    v_0(r) = c \cdot \frac{r^\eta}{r^\eta + K_{sig}^\eta}
\end{equation}

\subsection{Fit Results and Comparison}

The distributions were fitted to the 1701 reconstructed data points derived from the iterative rectification pipeline.

\begin{table}[h]
    \centering
    \begin{tabular}{|l|c|c|c|}
        \hline
        \textbf{Fit} & \textbf{Scale Radius ($K$)} & \textbf{Shape ($\eta$)} & \textbf{Fit Quality ($R^2$)} \\
        \hline
        \textbf{Exponential} & $\mathbf{4200}$ \textbf{Mpc} & N/A & $\mathbf{0.9929}$ \\
        Sigmoidal & $3029$ Mpc & $1.18$ & $0.9905$ \\
        \hline
    \end{tabular}
    \caption{\textbf{Comparison of analytical fits:} While both fits provide an excellent correlation ($R^2 > 0.99$), the Exponential fit achieves a higher $R^2$ with one fewer parameter.}
    \label{tab:potential_fit}
\end{table}

\subsection{The Implied Hubble Constant}

The "Local Hubble Constant" ($H_0$) is revealed, not as an accelerating expansion rate, but as the initial gradient of the potential well ($H_0 \equiv c/K$).

\subsubsection{Case A: The Exponential Prediction}
For the exponential fit, the Taylor expansion near the origin is linear ($v \approx c \cdot r/K$). Using the converged scale radius $K_{exp} \approx 4200.4$ Mpc:
\begin{equation}
    H_0^{(Exp)} = \frac{c}{K_{exp}} \approx \frac{299,792.5 \text{ km/s}}{4200.4 \text{ Mpc}} \approx \mathbf{71.4 \text{ km/s/Mpc}}
\end{equation}
This value arises directly from the fit and is in excellent agreement with local measurements (SH0ES $H_0 \approx 73$ km/s/Mpc) and Planck measurements ($H_0 \approx 67.4$ km/s/Mpc), effectively bridging the gap between them.

\subsubsection{Case B: The Sigmoidal Prediction}
For the Sigmoidal fit, the shape parameter $\eta \approx 1.18$ suppresses the velocity near the origin ($v \propto r^{1.18}$). To compensate for this slow start, the distribution forces the scale radius $K_{sig}$ to be significantly smaller ($3028.6$ Mpc). This implies a much steeper underlying scale factor:
\begin{equation}
    H_0^{(sig)} = \frac{c}{K_{sig}} \approx \frac{299,792.5 \text{ km/s}}{3028.6 \text{ Mpc}} \approx \mathbf{99.0 \text{ km/s/Mpc}}
\end{equation}
While the \textit{effective} slope in the observable range matches the data, the fundamental parameter $c/K$ diverges significantly from accepted values.

\subsection{Conclusion}

The analysis favors the \textbf{Exponential Fit} on two grounds:
\begin{enumerate}
    \item \textbf{Statistical Parsimony:} It achieves a higher $R^2$ score ($0.9929$ vs.\ $0.9905$) with a single free parameter rather than two.
    \item \textbf{Physical Consistency:} It recovers a local Hubble Constant ($71.4$ km/s/Mpc) consistent with observation, whereas the Sigmoidal model ($H_0 \approx 99.0$ km/s/Mpc) requires an anomalously steep scale factor to compensate for its non-linear core.
\end{enumerate}

% ----------------------------------------------------------------------
% Appendix C: Statistical Robustness Verification
% ----------------------------------------------------------------------
\section{Appendix C: Statistical Robustness Verification}
\label{app:robustness}

To ensure that the derived scale radius $K$ is a physical property of the cosmological potential and not an artifact of the inversion technique (Appendix~\ref{app:data_reconstruction}), we performed a secondary "Forward Modeling" analysis on the full Pantheon+ dataset ($N=1701$).

\subsection{Methodological Comparison}

We contrasted two distinct statistical approaches to determine the optimal potential parameters:

\begin{enumerate}
    \item \textbf{Method 1: Algebraic Inversion (Derived).} 
    The observational coordinates $(z, \mu)$ were transformed, point-by-point, into physical phase space $(r_{true}, v_0)$ using the relations derived in Section~\ref{sec:results_discussion} and the iterative rectification pipeline. We then fit the potential function $v_0(r)$ to this kinematic cloud to minimize the velocity residuals.
    
    \item \textbf{Method 2: Forward Optimization (Robustness).} 
    Luminosity distance was defined as $d_L(z; K) = (1+z)r_{opt}(z; K)$ directly in magnitude space using the exact closed-form integral $r_{opt}(z; K) = K[z/2 + \frac{1}{4}\ln(2z+1)]$. A non-linear optimization algorithm (Nelder-Mead) was then used to find the scale radius $K$ that minimizes the uncertainty-weighted $\chi^2$ statistic using the Pantheon+ diagonal uncertainties ($\sigma_{\mu,i} \equiv \texttt{MU\_SH0ES\_ERR\_DIAG}$):
    \begin{equation}
        \chi^2(K) = \sum_{i=1}^{N} \left( \frac{\mu_{obs, i} - \mu_{model}(z_i, K)}{\sigma_{\mu, i}} \right)^2
    \end{equation}
\end{enumerate}

\subsection{Results and $\Lambda$CDM Benchmark Comparison}

The two methods converged on the same physical scale radius for the Exponential potential to within $0.42\%$ ($K_{inv} = 4200.4$ Mpc vs.\ $K_{fwd} = 4182.7$ Mpc), demonstrating that the potential profile is the global minimum of the dataset regardless of the coordinate space used. 

Crucially, Table~\ref{tab:robustness_results} also demonstrates that the exact refractive conversion $r_{true} \to r_{opt} = \int_0^{r_{true}} (1 - v_0^2/c^2)^{-1}\,dr = K[z/2 + \frac{1}{4}\ln(2z+1)]$ is not a coordinate relabeling of pure coasting FLRW ($n = 1$, $d_L = (1+z)K\ln(1+z)$). Evaluating the exact refractive integral reduces $\chi^2$ by $\Delta\chi^2 = -64.15$ relative to pure coasting FLRW (from $839.29$ down to $775.14$) at zero additional parameters, bringing the 1-parameter refractive model within $\Delta\chi^2 = +25.93$ of 1-parameter Planck-prior $\Lambda$CDM ($\Omega_m = 0.315$, $\chi^2 = 749.21$) with $\Omega_\Lambda = 0$.

\begin{table}[h]
    \centering
    \renewcommand{\arraystretch}{1.25} 
    \resizebox{\textwidth}{!}{%
    \begin{tabular}{lccccc}
        \toprule
        \textbf{Model ($N = 1701$)} & \textbf{Params ($k$)} & \textbf{Inverse Fit} & \textbf{Forward Fit ($\min \chi^2$)} & \textbf{$\chi^2$ ($\chi^2/\text{dof}$)} & \textbf{BIC} \\
        \midrule
        \textbf{Refractive Exponential} & $1$ & $K = 4200.4$ Mpc & $K = 4182.7$ Mpc ($\mathbf{0.42\%}$) & $\mathbf{775.14}$ ($\mathbf{0.456}$) & $\mathbf{782.58}$ \\
        \textbf{Refractive Sigmoidal} & $2$ & $K = 3028.6, \eta = 1.18$ & $K = 3379.0, \eta = 1.06$ & $1159.32$ ($0.682$) & $1174.19$ \\
        \midrule
        \textbf{Pure Coasting FLRW ($n = 1$)} & $1$ & --- & $K = 4247.1$ Mpc & $839.29$ ($0.494$) & $846.73$ \\
        \textbf{Un-lensed Linear Hubble Law} & $1$ & --- & $H_0 = 64.56$ km/s/Mpc & $2458.54$ ($1.446$) & $2465.97$ \\
        \textbf{Flat $\Lambda$CDM (Planck $\Omega_m = 0.315$)} & $1$ & --- & $H_0 = 73.36$ km/s/Mpc & $749.21$ ($0.441$) & $756.65$ \\
        \textbf{Flat $\Lambda$CDM (2-Parameter Free)} & $2$ & --- & $H_0 = 72.97, \Omega_m = 0.351$ & $745.40$ ($0.439$) & $760.28$ \\
        \bottomrule
    \end{tabular}%
    }
    \caption{\textbf{Statistical Robustness and Pantheon+ $\chi^2$ Model Comparison ($N=1701$).} Comparison of inverse velocity reconstruction ($r_{true}$) against forward $\chi^2$ minimization of weighted magnitude residuals under the exact metric refractive index $n = (1 - v_0^2/c^2)^{-1}$, alongside standard cosmological baselines.}
    \label{tab:robustness_results}
\end{table}

\begin{figure}[H]
    \centering
    \includegraphics[width=\linewidth]{robustness_comparison_final.png}
    \caption{\textbf{Robustness of the Cosmological Potential.} 
    \textbf{Panel A (Top):} The derived Exponential potential (Blue Solid, $K=4200$ Mpc) is overlaid with the Forward-optimized model (Cyan Dashed, $K=4183$ Mpc). The curves are virtually indistinguishable ($0.42\%$ agreement), confirming that minimizing magnitude errors yields the same physical geometry as inverting velocity data.
    \textbf{Panel B (Bottom):} The Sigmoidal model comparison between the Inverse (Red Solid, $K=3029, \eta=1.18$) and Forward (Orange Dashed, $K=3379, \eta=1.06$) methods.}
    \label{fig:robustness}
\end{figure}

% ----------------------------------------------------------------------
% Appendix D: The High-Redshift Chronometer
% ----------------------------------------------------------------------
\section{Appendix D: The High-Redshift Chronometer}
\label{app:time_dilation}

A critical tension in contemporary cosmology is the discovery of massive, mature galaxies at extreme redshifts ($z \approx 14$)~\cite{Labbe2023}. Under standard matter-dominated expansion, the proper age of the universe at this epoch leaves insufficient time for the formation of the observed stellar populations.

The compound relativistic time-dilation factor between a free-fall source at radial depth $r$ and the observer at $r=0$ is governed by the product of gravitational time dilation and kinematic Doppler dilation derived in Appendix~\ref{app:data_reconstruction}:
\begin{equation}
    \frac{d\tau}{dt_0} = \frac{1}{1+z(r)} = 1 - \frac{v_0(r)}{c}
\end{equation}

Along the radial potential well, the differential proper time remaining relative to the asymptotic horizon ($r \to \infty$, where $v_0 \to c$ and $d\tau/dt_0 \to 0$) is $d\tau = \frac{1}{c}(1 - v_0(r)/c)\,dr$. Substituting the Exponential velocity profile $v_0(r) = c(1 - e^{-r/K})$ gives $1 - v_0(r)/c = e^{-r/K}$. Integrating this proper-time density from a source at $r_{true}(z) = K\ln(1+z)$ out to the horizon ($r \to \infty$) yields a finite, closed-form proper chronometer anchored directly to the spatial scale radius $K$:
\begin{equation}
    \tau(z) = \int_{r_{true}(z)}^{\infty} \frac{1}{c}\left(1 - \frac{v_0(r')}{c}\right) dr' = \int_{K\ln(1+z)}^{\infty} \frac{e^{-r'/K}}{c}\,dr' = \frac{K}{c}e^{-\ln(1+z)} = \frac{K/c}{1+z}
\end{equation}

At the observer's origin ($z = 0, r = 0$), the present proper age evaluates to $\tau_0 = K/c$. Integrating the complementary interval from $r = 0$ to $r_{true}(z)$ yields the chronological look-back time $\tau_{\text{lookback}}(z)$:
\begin{equation}
    \tau_{\text{lookback}}(z) = \int_0^{r_{true}(z)} \frac{e^{-r'/K}}{c}\,dr' = \tau_0\left(1 - \frac{1}{1+z}\right) = \tau_0\frac{z}{1+z}
\end{equation}
This satisfies the exact chronological sum rule $\tau(z) + \tau_{\text{lookback}}(z) = \tau_0$ for all $z \ge 0$, distinguishing the bounded chronological look-back time ($\tau_{\text{lookback}} < \tau_0$) from the spatial optical lever arm $r_{opt}(z) = \int_0^{r_{true}(z)} n(r')\,dr'$ that governs inverse-square flux dilution.

\subsection{Empirical Evaluation (The Early Galactic Boundary)}

Evaluating this potential chronometer at extreme high-redshift observations ($z = 14$ and $z = 20$) isolates the proper cosmological age at the epoch of emission. 

Using our empirically derived Exponential scale radius ($K \approx 4200.4 \text{ Mpc}$), the baseline proper age evaluates to $\tau_0 = K/c \approx 13.70 \text{ Gyr}$ (with a chronological look-back time to $z=14$ of $\tau_{\text{lookback}}(14) = \frac{14}{15}\tau_0 \approx 12.79 \text{ Gyr}$). The proper evolutionary age at these extreme horizons evaluates to:
\begin{align}
    \tau(14) &= \frac{13.70 \text{ Gyr}}{1 + 14} \approx 913 \text{ Myr} \\
    \tau(20) &= \frac{13.70 \text{ Gyr}}{1 + 20} \approx 652 \text{ Myr}
\end{align}

The true proper evolutionary window for the galaxy is the difference between these two epochs ($\Delta t_{evol} = \tau(14) - \tau(20)$). Table~\ref{tab:time_window} demonstrates that this continuous temporal scaling provides significantly more proper time for galactic evolution than the standard model.

\begin{table}[h]
    \centering
    \begin{tabular}{|l|c|c|c|}
        \hline
        \textbf{Model Baseline} & \textbf{Metric Age ($z=14$)} & \textbf{Proper Window ($z=20 \to 14$)} & \textbf{Factor vs $\Lambda$CDM} \\
        \hline
        $\Lambda$CDM & $\approx 285$ Myr & $\approx 137$ Myr & $1.0\times$ \\
        \textbf{Exponential Fit} & $\mathbf{913}$ \textbf{Myr} & $\mathbf{261}$ \textbf{Myr} & $\mathbf{\approx 1.9\times}$ \\
        \hline
    \end{tabular}
    \caption{\textbf{The proper time available for galactic evolution.} By integrating the compound relativistic time-dilation factor across the Exponential potential well, the metric provides nearly double the proper chronological window for stellar evolution compared to standard $\Lambda$CDM.}
    \label{tab:time_window}
\end{table}

This affords the required time for galactic evolution to match the observed, massive galaxies at extreme optical horizons.\footnote{\textbf{Note on the Sigmoidal Distribution:} For the Sigmoidal profile $v_0(r) = c (r/K)^\eta / [1 + (r/K)^\eta]$ ($\eta \approx 1.18$), the horizon residual $1 - v_0(r)/c = [1 + (r/K)^\eta]^{-1}$ decays algebraically rather than exponentially. Consequently, its tail integral fails to saturate on a physical horizon timescale ($\tau(14) \approx 37.0\text{ Gyr}$), further favoring the Exponential potential.}

% ----------------------------------------------------------------------
% Appendix E: The Origin of the Dark Energy Signal
% ----------------------------------------------------------------------
\section{Appendix E: The Origin of the Dark Energy Signal}
\label{app:dark_energy_origin}

Dark energy originated as a solution to the discrepancy between the observed brightness of Type Ia supernovae and the predictions of a linear expansion model. In this appendix, we derive the source of this discrepancy using the two datasets generated in Appendix~\ref{app:data_reconstruction}.

\subsection{The Two-Metric Problem}

In Appendix~\ref{app:data_reconstruction}, we separated the observational data into two distinct coordinate systems:
\begin{enumerate}
    \item \textbf{Dataset 1 (Apparent Space $r_{opt}$):} The path-distance light \textit{appears} to have traveled, based on the observed flux. This includes all optical distortions caused by the vacuum potential.
    \item \textbf{Dataset 2 (Physical Space $r_{true}$):} The displacement the light \textit{actually} traveled, obtained by rectifying (de-lensing) the path through the potential gradient.
\end{enumerate}

"Dark Energy" is potentially revealed to be an observational signal where high-redshift supernovae appear dimmer (and thus farther away) than a linear matter-only model predicts. We can now quantify this anomaly $\Delta \mu$ as the divergence between our two datasets.

\subsection{Deriving the Signal}

The magnitude $\mu$ is logarithmically related to the luminosity distance $d_L$:
\begin{equation}
    \mu = 5 \log_{10}(d_L) + 25
\end{equation}

By Etherington's reciprocity relation, the luminosity distance is related to the coordinate distance by the redshift factor $(1+z)$:
\begin{itemize}
    \item \textbf{Observation (Apparent):} $\mu_{obs}$ corresponds to the apparent coordinate $r_{opt}$.
    \item \textbf{Expectation (Un-lensed):} $\mu_{lin}$ corresponds to the physical coordinate $r_{true}$.
\end{itemize}

The magnitude anomaly (the signal usually attributed to accelerated expansion) is the difference between what we observe and what we expect geometrically:
\begin{equation}
    \Delta \mu_{DE} = \mu_{obs} - \mu_{lin}
\end{equation}

Substituting the coordinate definitions:
\begin{equation}
    \Delta \mu_{DE} = \left[ 5 \log_{10}(r_{opt}(1+z)) + 25 \right] - \left[ 5 \log_{10}(r_{true}(1+z)) + 25 \right]
\end{equation}

The redshift factors $(1+z)$ and constants cancel out perfectly, isolating the refractive origin of the signal:
\begin{equation}
    \Delta \mu_{DE} = 5 \log_{10} \left( \frac{r_{opt}}{r_{true}} \right)
\end{equation}

\begin{figure}[H]
    \centering
    \includegraphics[width=0.85\textwidth]{hubble_residuals_bridge.png}
    \caption{\textbf{Resolution of the Magnitude Anomaly.}
    \textit{Top Panel:} The standard "Dark Energy" signal (grey points), calculated as the deviation from a linear Hubble Law ($\mu_{obs} - \mu_{linear}$). The solid blue line represents the theoretical refractive prediction derived from the Exponential Potential ($K \approx 4200$ Mpc) under the exact metric index $n = (1 - v_0^2/c^2)^{-1}$.
    \textit{Bottom Panel:} The residuals after applying the refractive correction (blue points). The signature vanishes ($\Delta \mu \approx 0$) across all redshifts, demonstrating that the anomaly is accounted for by the optical path lengthening.}
    \label{fig:hubble_bridge}
\end{figure}

\subsection{Empirical Verification}

We defined the exact refractive index of the vacuum potential as $n(r) = (1 - v_0^2/c^2)^{-1}$. The apparent distance $r_{opt}$ is the integral of this index along the line of sight:
\begin{equation}
    r_{opt} = \int_0^{r_{true}} n(r') \, dr' = \int_0^{r_{true}} \left(1 - \frac{v_0^2(r')}{c^2}\right)^{-1} dr' > r_{true}
\end{equation}

Because $n(r) > 1$ everywhere along the path (due to the non-zero potential), the optical path $r_{opt}$ is strictly longer than the geometric path $r_{true}$:
\begin{itemize}
    \item As redshift increases, the photon traverses a longer path through the potential.
    \item The ratio $r_{opt} / r_{true}$ grows continuously.
    \item The magnitude anomaly $\Delta \mu_{DE}$ increases positively.
\end{itemize}

Figure~\ref{fig:hubble_bridge} confirms this relationship empirically. When we subtract the refractive distortion ($5 \log_{10}(r_{opt}/r_{true})$) from the observed data, the "acceleration" curve flattens completely. This demonstrates that Dark Energy in the supernova Hubble diagram can be understood as an optical property of the vacuum potential.

% ----------------------------------------------------------------------
% Appendix F: The Metric of the Vacuum Potential
% ----------------------------------------------------------------------
\section{Appendix F: The Metric of the Vacuum Potential}
\label{app:vacuum_metric}

In the preceding appendices, we reconstructed the velocity potential $v_0(r)$ and evaluated its exact refractive index $n(r) = (1 - v_0^2/c^2)^{-1}$. In this appendix, we summarize the embedding of this potential into the reciprocal compliance metric of Section~\ref{sec:gravitational_gradient}.

\subsection{The Static Metric of a Refractive Vacuum}

Refractive Cosmology models the past light cone via a static, spherically symmetric metric governed by the global potential $\Phi(r) = -\frac{1}{2}v_0^2(r)$. As introduced in Eq.~\ref{eq:v0_metric}, the \textbf{Vacuum Potential Metric} is:
\begin{equation}
    ds^2 = -\left(1 + \frac{2\Phi(r)}{c^2}\right) c^2 dt^2 + \left(1 + \frac{2\Phi(r)}{c^2}\right)^{-1} dr^2 + r^2 d\Omega^2
\end{equation}
where $1 + 2\Phi(r)/c^2 = 1 - v_0^2(r)/c^2$.

\subsection{Unifying Matter and Light}

\subsubsection{Matter Geodesics (The Potential Flow)}
Galaxies follow timelike geodesics ($u^\mu u_\mu = -1$). In the weak-field regime, the radial equation of motion is:
\begin{equation}
    \frac{d^2r}{d\tau^2} = -\Gamma^r_{tt} \left(\frac{dt}{d\tau}\right)^2 \approx -\partial_r \Phi
\end{equation}
The recession velocity field is identified as the free-fall kinematic response of matter to the global potential gradient ($v_0^2 = -2\Phi$).

\subsubsection{Null Geodesics (The Refractive Index)}
Light follows null geodesics ($ds^2 = 0$). For a radial ray ($d\Omega=0$), Eq.~\ref{eq:v0_metric} simplifies to:
\begin{equation}
    0 = -\left(1 - \frac{v_0^2}{c^2}\right) c^2 dt^2 + \left(1 - \frac{v_0^2}{c^2}\right)^{-1} dr^2
\end{equation}
Rearranging for the coordinate speed of light $v_c(r) = dr/dt$ yields:
\begin{equation}
    \frac{dr}{dt} = c \left( 1 - \frac{v_0^2}{c^2} \right)
\end{equation}
This recovers the exact refractive index $n(r) \equiv c/v_c = (1 - v_0^2/c^2)^{-1}$. Substituting the Exponential velocity profile $v_0(r) = c(1 - e^{-r/K})$ and $r_{true}(z) = K\ln(1+z)$ yields the closed-form optical path integral $r_{opt}(z) = K[z/2 + \frac{1}{4}\ln(2z+1)]$ utilized throughout Section~\ref{sec:results_discussion}, Appendix~\ref{app:data_reconstruction}, Appendix~\ref{app:robustness}, and Appendix~\ref{app:dark_energy_origin}.

\subsection{Effective Refractive Analogues of the Friedmann Equations}
By mapping the radial refractive gradient to effective analogues of the standard cosmological parameters, we define two phenomenological identifications:

\begin{enumerate}
    \item \textbf{The Gradient Identification (Effective Energy Density):}
    The square of the radial refractive gradient serves as the optical analogue to the Friedmann expansion rate $H^2$:
    \begin{equation}
        \left( c \frac{\partial_r n}{n^2} \right)^2 = \frac{8\pi G}{3} \rho_{\text{eff}} - \frac{c^2}{K^2}
    \end{equation}

    \item \textbf{The Ray Identification (Effective Acceleration):}
    The apparent longitudinal acceleration is governed by the Laplacian of the inverse refractive index:
    \begin{equation}
        \frac{c^2}{n} \nabla^2 \left( \frac{1}{n} \right) = -4\pi G (\rho + 3p_{\text{eff}})
    \end{equation}
\end{enumerate}

\subsection{Application: Velocity-Slope vs. BAO Observables}
For the Exponential Potential model ($v_0(r) = c(1-e^{-r/K})$) with $r_{true}(z) = K \ln(1+z)$, the model yields two distinct, testable functions of redshift (Section~\ref{sec:Hubble_Tension} and Section~\ref{sec:bao_delensing}):
\begin{align}
    H_{slope}(z) &\equiv \frac{dv_0}{dr} = \frac{c}{K(1+z)} = \frac{H_0}{1+z}, \\
    D_{H,true}(z) &\equiv \frac{dr_{true}}{dz} = \frac{K}{1+z} \implies H_{BAO}(z) \equiv \frac{c}{D_{H,true}(z)} = H_0(1+z).
\end{align}
This creates specific testable predictions:
\begin{itemize}
    \item \textbf{Low-$z$ ($z \ll 1$):} Both $H_{slope}(z)$ and $H_{BAO}(z)$ converge to the identical local intercept $H_0 = c/K \approx 71.4$ km/s/Mpc.
    \item \textbf{High-$z$ ($z \sim 1\text{--}2.3$):} The velocity-distance slope $H_{slope}(z)$ flattens as $H_0/(1+z)$, resolving the Hubble Tension between local tangent measurements and high-redshift baseline fits, while $D_{H,true}(z) = K/(1+z)$ establishes the pure-vacuum baseline onto which baryonic matter perturbations ($\Omega_b(1+z)^3$) couple at intermediate BAO redshifts (Section~\ref{sec:bao_delensing}).
\end{itemize}

\begin{table}[h]
\centering
\begin{tabular}{l|l}
\toprule
FLRW Cosmology & Refractive Cosmology \\
\midrule
Scale factor $a(t)$ & Refractive index $n(r)$ \\
Hubble $H = \dot{a}/a$ & Gradient $c \partial_r n / n^2$ \\
Curvature $k/a^2$ & Basin scale $1/K^2$ \\
Dark energy $\Lambda$ & Horizon saturation $v_0 \to c$ \\
\bottomrule
\end{tabular}
\caption{Translation between FLRW and refractive descriptions.}
\end{table}

\subsection{Conclusion: The Universe as a Gradient-Index Lens}
The apparent acceleration of the cosmos is modeled as the longitudinal optical delay of a global Gradient-Index (GRIN) lens embedded in General Relativity, arising from the saturation of the potential toward the cosmic horizon ($v_0 \to c$) without requiring negative pressure or an exotic dark energy fluid.

% ----------------------------------------------------------------------
% Appendix G: The Cosmic Schlieren Test
% ----------------------------------------------------------------------
\section{Appendix G: The Cosmic Schlieren Test - A Proposal for Independent Verification}
\label{app:cosmic_schlieren}

The central claim of this work is that the vacuum possesses a global refractive index gradient $\nabla n(r)$ governed by the scalar velocity field $v_0(r)$. While we have validated this claim using Supernova distance moduli (longitudinal lensing), a definitive independent test can be performed using \textbf{CMB Lensing Reconstruction}.

\subsection{The Universe as a Schlieren Interferometer}
A classic Schlieren optical system visualizes fluid dynamics by detecting the deflection of collimated light caused by gradients in the refractive index~\cite{Toepler1864}. The universe effectively constitutes a natural Schlieren setup:
\begin{itemize}
    \item \textbf{Backlight:} The Cosmic Microwave Background (CMB) provides a uniform, collimated source screen at $z \approx 1100$.
    \item \textbf{Medium:} The gravitational potential $v_0(r)$ acts as the refractive fluid with index $n(r)$.
    \item \textbf{Signal:} The deflection field $\mathbf{d}(\hat{\mathbf{n}})$ acts as the Schlieren image, revealing the gradients $\nabla n$.
\end{itemize}

\subsection{The Predicted Signal: Global Convergence}
Standard cosmological analysis assumes that the universe is statistically isotropic and homogeneous on large scales. Consequently, lensing reconstruction pipelines typically subtract the monopole ($\ell=0$) and dipole ($\ell=1$) modes of the lensing potential $\phi$, interpreting them as systematic noise or motion.

Refractive lensing, however, predicts that the background potential is not flat. The scalar field $v_0(r)$ possesses a non-zero Laplacian $\nabla^2 v_0$, which generates a \textbf{Global Convergence} term ($\kappa_0$) in the lensing map:
\begin{equation}
    \kappa_{global} = \frac{1}{2} \nabla^2 \phi_{eff} \propto \nabla^2 n(r) \neq 0.
\end{equation}
Unlike random Dark Matter fluctuations which average to zero ($\langle \kappa \rangle \approx 0$), this global convergence creates a systematic isotropic magnification (or demagnification) of the CMB power spectrum.

\subsection{Connection to the $A_L$ Anomaly}
This structural feature may provide a geometric mechanism for the signal reported in the Planck 2018 results. Planck reports a persistent "Lensing Amplitude Anomaly," where the smoothing of the acoustic peaks implies an excess lensing amplitude parameter of $A_L \approx 1.18$~\cite{Planck2018}.

We hypothesize that this excess may be the Schlieren signature of CMB photons traversing the uncorrected macroscopic scalar velocity field $v_0(r)$, rather than an instrumental error or statistical fluke.

\subsection{Proposed Test Methodology}
We propose a direct test using the open-access Planck Legacy Archive (PLA) data, effectively treating the CMB analysis as a Global Schlieren Interferometer. By tracking the directional structural variance of the background potential to the $z=1100$ horizon, we can distinguish between random local clustering and coherent global geometry.

\begin{enumerate}
    \item \textbf{Unfiltered Reconstruction:} Standard pipelines filter out low-multipole ($\ell < 10$) modes, assuming them to be systematic noise. We propose performing a lensing reconstruction on the raw temperature maps \textit{without} this subtraction to preserve the signal of the global gradient $\nabla n$.
    
    \item \textbf{Directional Variance Mapping:} If the observer is slightly offset from the center of the potential well, the global gradient $\nabla n$ will manifest as a coherent directional bias in the lensing potential. This should appear as a strong dipole in the lensing map aligned with the matter distribution, distinct from the kinematic CMB dipole.
    
    \item \textbf{Radial Cross-Correlation:} Cross-correlate the lensing convergence map $\kappa(\hat{\mathbf{n}})$ with the Supernova magnitude residuals. A positive correlation would confirm that the same potential structure responsible for the "Dark Energy" dimming at low redshift is responsible for the excess lensing variance at high redshift.

    \item \textbf{Tomographic Focusing:} Because the refractive index $n(r)$ is depth-dependent, the universe acts as a Gradient-Index (GRIN) lens. By applying tomographic techniques to sources at varying redshifts (e.g., combining CMB, CIB, and Galaxy Weak Lensing), we can effectively "focus" the Schlieren apparatus on specific radial shells. This would allow for the 3D reconstruction of the vacuum potential, revealing not just the global gradient, but the local refractive structure of cosmic filaments and voids.
\end{enumerate}

While a full execution of this Tomographic Schlieren analysis remains beyond the scope of this paper, the methodology outlined here provides a clear path for independent verification using existing public data.

% ======================================================================
% BIBLIOGRAPHY
% ======================================================================
\begin{thebibliography}{99}

\bibitem{Adame2024}
Adame, A. G., et al. (DESI Collaboration) (2024).
``DESI 2024 VI: Cosmological Constraints from the Measurements of Baryon Acoustic Oscillations.''
\textit{arXiv preprint}, arXiv:2404.03002.

\bibitem{Alam2021}
Alam, S., et al. (2021).
``Completed SDSS-IV extended Baryon Oscillation Spectroscopic Survey: Cosmological implications from two decades of spectroscopic surveys at the Apache Point Observatory.''
\textit{Physical Review D}, \textbf{103}(8), 083533.

\bibitem{RefractivePackage2026}
Bair-Mosheh, B. T. (2026).
``Refractive Cosmology Reproduction Package (v2.1).''
\textit{Code and Data Artifact}.
Available as Supplementary Material to this article.

\bibitem{Bartelmann2001}
M.~Bartelmann and P.~Schneider, 
\textit{Weak gravitational lensing}, 
Phys. Rep. \textbf{340}, 291 (2001).

\bibitem{Beutler2011}
Beutler, F., et al. (2011). 
``The 6dF Galaxy Survey: baryon acoustic oscillations and the local Hubble constant.'' 
\textit{Monthly Notices of the Royal Astronomical Society}, \textbf{416}(4), 3017--3032.

\bibitem{BornWolf}
Born, M., \& Wolf, E. (1999).
\textit{Principles of Optics: Electromagnetic Theory of Propagation, Interference and Diffraction of Light}.
Cambridge University Press.

\bibitem{PantheonPlus}
Brout, D., et al. (2022).
``The Pantheon+ Analysis: Cosmological Constraints.''
\textit{The Astrophysical Journal}, \textbf{938}(2), 110.

\bibitem{Buchert2000}
Buchert, T. (2000).
``On average properties of inhomogeneous cosmologies.''
\textit{General Relativity and Gravitation}, \textbf{32}, 105--125.

\bibitem{Eddington1920}
Eddington, A. S. (1920).
\textit{Space, Time and Gravitation: An Outline of the General Relativity Theory}.
Cambridge University Press.

\bibitem{Einstein1916}
Einstein, A. (1916).
``Die Grundlage der allgemeinen Relativit\"atstheorie.''
\textit{Annalen der Physik}, \textbf{354}(7), 769--822.

\bibitem{Gordon1923}
Gordon, W. (1923).
``Zur Lichtfortpflanzung nach der Relativit\"atstheorie.''
\textit{Annalen der Physik}, \textbf{377}(22), 421--456.

\bibitem{Labbe2023}
Labb\'e, I., et al. (2023).
``A population of red candidate massive galaxies $\sim$600 Myr after the Big Bang.''
\textit{Nature}, \textbf{616}, 266--269.

\bibitem{Minkowski1909}
Minkowski, H. (1909). 
``Raum und Zeit.'' \textit{Physikalische Zeitschrift}, \textbf{10}, 104-111. 
(Address delivered at the 80th Assembly of German Natural Scientists and Physicians, Cologne, 1908).

\bibitem{Planck2018}
Planck Collaboration. (2020).
``Planck 2018 results. VI. Cosmological parameters.''
\textit{Astronomy \& Astrophysics}, \textbf{641}, A6.

\bibitem{Poisson2004}
Poisson, E. (2004). 
\textit{A Relativist's Toolkit: The Mathematics of Black-Hole Mechanics}. 
Cambridge University Press.

\bibitem{Poisson2014}
Poisson, E., \& Will, C. M. (2014). 
\textit{Gravity: Newtonian, Post-Newtonian, Relativistic}. 
Cambridge University Press.

\bibitem{Pound1960}
Pound, R. V., \& Rebka Jr, G. A. (1960). 
``Apparent Weight of Photons.''
\textit{Physical Review Letters}, \textbf{4}(7), 337.

\bibitem{SHOES2022}
Riess, A. G., et al. (2022).
``A Comprehensive Measurement of the Local Value of the Hubble Constant with 1 km/s/Mpc Uncertainty from the Hubble Space Telescope and the SH0ES Team.''
\textit{The Astrophysical Journal Letters}, \textbf{934}(1), L7.

\bibitem{Shapiro1964}
Shapiro, I. I. (1964).
``Fourth Test of General Relativity.''
\textit{Physical Review Letters}, \textbf{13}(26), 789--791.

\bibitem{Toepler1864}
Toepler, A. (1864). \textit{Beobachtungen nach einer neuen optischen Methode: Ein Beitrag zur Experimental-Physik}. Bonn: Max Cohen \& Sohn.

\bibitem{Soldner1801}
von S\"oldner, J. (1801).
``Ueber die Ablenkung eines Lichtstrals von seiner geradlinigen Bewegung...''
\textit{Berliner Astronomisches Jahrbuch}, 161--172.

\bibitem{Weinberg1976}
Weinberg, S. (1976).
``Apparent luminosities in a locally inhomogeneous universe.''
\textit{The Astrophysical Journal}, \textbf{208}, L1--L3.

\end{thebibliography}

\end{document}